\section{Conclusiones}\label{sec:conclusiones}

\subsection{Aportes del diseño}

El sistema demuestra que una memoria compartida distribuida puede ofrecer una interfaz de punteros normales sin renunciar a la coherencia ni a la tolerancia a fallos. La vista se proyecta sobre un \texttt{memfd} en una dirección común y se protege página a página con \texttt{mprotect}; la aplicación lee y escribe como si tuviera un bloque de memoria local, mientras el sistema resuelve cada fallo con el protocolo MSI. Esa decisión mantiene el modelo de programación simple y traslada la complejidad al núcleo, que es donde puede tratarse de forma uniforme para todas las aplicaciones.

El manejador de \texttt{SIGSEGV} es la pieza que hace posible la transparencia. El manejador no hace red, no toma locks y no reserva memoria: solo identifica la página, registra el fallo y duerme al hilo hasta que un despachador lo resuelve. Esa disciplina convierte una señal que normalmente termina el proceso en un mecanismo de demanda de páginas. El diseño separa el camino de la señal del camino de la red y evita que una espera de red dentro del manejador bloquee al proceso o introduzca dependencias circulares.

El directorio en el nodo custodio y el anillo de dispersión consistente eliminan la necesidad de un coordinador de propiedad. Cada página tiene un home que conoce sus copias, su dueño y su versión; el home se calcula de forma determinista a partir de la membresía y del UUID de cada nodo, ponderado por la memoria que aporta. Todos los nodos llegan al mismo resultado sin consultar a nadie, y cuando la membresía cambia solo se mueven las páginas de los tramos afectados. La combinación de directorio y anillo permite que la coherencia sobreviva a la caída de un custodio y que el contenido se transfiera al nuevo home con \texttt{PAGE\_MOVE}.

Los estados transitorios del protocolo MSI son otro aporte del diseño. Una transición no ocurre en el instante en que se solicita: la página pasa por \texttt{IS\_PEND}, \texttt{IM\_PEND} o \texttt{SM\_PEND} mientras la respuesta viaja por la red. Esos estados permiten describir la página durante la espera sin mantener un lock tomado, y los permisos de la vista reflejan lo que el nodo puede hacer en cada momento. El resultado es un protocolo que espera sin bloquear el resto del sistema y que puede agrupar solicitudes concurrentes sobre la misma página.

La redundancia con paridad completa el modelo con una garantía de disponibilidad. Las franjas agrupan $K$ páginas de datos y una de paridad, con la paridad en un custodio distinto de los datos. El delta de paridad se calcula al persistir, no en cada almacenamiento, y la reconstrucción por XOR recupera una pieza perdida a partir de las supervivientes. La validación confirma que la caída de un nodo no destruye el contenido de las franjas recuperables y que el sistema puede seguir leyendo después de la baja. La paridad no reemplaza la coherencia; la complementa para que una página no se pierda cuando su custodio desaparece.

La ejecución distribuida y los reportes HTML amplían el sistema más allá de la memoria compartida. Las particiones se describen con direcciones de la vista, los kernels se registran por proceso y los resultados se agregan con CRC32; los reportes convierten el estado interno en una tabla legible sin detener el clúster. Esa combinación de API, mensajes y observabilidad hace que el sistema pueda usarse como plataforma de aplicaciones y no solo como una demostración de coherencia.

\subsection{Lecciones de implementación}

La primera lección es que las restricciones del manejador de señales son parte del diseño y no un detalle de codificación. El manejador solo puede usar llamadas seguras para señales, y cualquier intento de hacer red, reservar memoria o tomar un mutex dentro de él introduce un riesgo de interbloqueo o de corrupción. La solución consiste en separar el registro del fallo de su resolución: el manejador deposita el fallo en una estructura por hilo y despierta a un despachador, que es un hilo normal y puede esperar. Esa separación es la que permite que el camino del fallo siempre termine, incluso cuando la red es lenta o un par no responde.

La segunda lección es que los trabajadores no deben esperar respuestas. Un trabajador que atiende una solicitud inicia la transacción, envía las invalidaciones necesarias y vuelve al pool; la respuesta final la cierra otro trabajador cuando llega el último acuse. Si los trabajadores esperaran, dos solicitudes cruzadas podrían formar un ciclo y bloquear el nodo. La regla de que solo los despachadores esperan, y de que lo hacen sobre respuestas que producen trabajadores que no esperan, mantiene el grafo de espera acíclico. Es una regla sencilla de enunciar y fácil de romper sin darse cuenta al agregar un camino nuevo.

El orden de \texttt{mprotect} también enseña una lección concreta. Al degradar una copia propia, el sistema cierra la escritura primero, copia el contenido después y solo entonces marca la página como compartida. Si copiara antes de cerrar la escritura, un hilo de la aplicación podría modificar la página entre la copia y el cambio de permiso, y el contenido entregado al solicitante no incluiría esa modificación. La secuencia parece un detalle menor, pero es la diferencia entre una copia consistente y una copia que puede quedar desactualizada en una ventana de pocos microsegundos.

La compuerta de \texttt{tseq} y el mapa de acuses resuelven dos problemas que no se ven en una prueba feliz. Una respuesta duplicada o reordenada puede llegar después de una versión más nueva; si se instalara, la página retrocedería. El contador de secuencia por página permite descartar la respuesta vieja antes de tocar el contenido. De forma parecida, una escritura que invalida varias copias necesita saber que todas llegaron, no solo la última; el mapa de acuses por transacción cuenta cada destinatario una vez y no se deja engañar por un duplicado. Esos dos mecanismos convierten el protocolo en algo que resiste la red real y no solo el orden ideal de una simulación.

La heterogeneidad de memoria obliga a separar lo compartido de lo local. El hash de configuración incluye las claves que deben coincidir y deja fuera \texttt{NODE\_MEM}, los plazos y los hilos. Esa separación permite que un equipo con más memoria se una al mismo clúster sin obligar a los demás a declarar una capacidad que no tienen. La validación física confirma que la regla funciona: cuatro equipos con 256, 128, 128 y 128 MiB comparten la misma vista y producen el mismo resultado. La configuración deja de ser un archivo idéntico para todos y pasa a ser un contrato común con parámetros locales.

Las pruebas con procesos reales descubren problemas que una prueba unitaria no ve. Los escenarios de extremo a extremo arrancan nodos, eligen puertos, esperan marcadores y comparan resultados; esa combinación ejercita el descubrimiento, la unión, la coherencia y el cierre en el mismo orden en que lo haría un usuario. La prueba de paridad, en particular, elige una víctima y páginas reconstruibles a partir del directorio real, de modo que no depende de una colocación supuesta. Esa forma de probar es más lenta que una simulación, pero es la que confirma que el sistema funciona cuando la red, los puertos y los procesos se comportan como en el laboratorio.

\subsection{Extensiones posibles}

Una ampliación deseable consiste en dar a la vista un almacenamiento persistente. El respaldo y la paridad viven en memoria; llevarlos a un archivo por nodo o a un dispositivo de memoria persistente permitiría conservar las páginas entre arranques y convertir \texttt{dsm\_sync} en un punto de recuperación. La misma estructura de marcos y versiones serviría para escribir un registro de cambios que reconstruya el estado después de una parada.

El reequilibrio dinámico de páginas es otra extensión natural. Cuando entra o sale un nodo, el anillo cambia y algunas páginas cambian de custodio; el sistema ya transfiere el contenido con \texttt{PAGE\_MOVE}. Un reequilibrio gradual permitiría mover las páginas en segundo plano, con un límite de transferencias por intervalo, para que una membresía grande no produzca una ráfaga que sature la red. El mismo mecanismo podría repartir las páginas de una franja entre nodos distintos y acercar la tolerancia de una pieza por franja a la tolerancia de un nodo completo.

La tolerancia a particiones de red exige un paso adicional sobre el modelo de fallos de parada. Un nodo que deja de responder no siempre está muerto: puede estar aislado y seguir ejecutando operaciones. Un esquema de quórum, arrendamientos o épocas con confirmación mutua permitiría distinguir una partición de una caída y evitar que dos mitades del clúster se declaren dueñas de la misma página. Esa extensión cambia las reglas de membresía y de directorio, pero conserva la idea de que cada página tiene un custodio y una versión.

Las transacciones multiescritor abren otra línea. La coherencia actual garantiza una página con un solo escritor o varios lectores, no una operación atómica sobre varias páginas. Una capa de transacciones con versiones, detección de conflictos y confirmación en dos fases permitiría actualizar estructuras enlazadas o índices sin que el programador tenga que construir su propio protocolo con \texttt{dsm\_lock}. La misma capa podría ofrecer memoria transaccional por software sobre la vista y reutilizar los cerrojos distribuidos como base.

El transporte también admite mejoras. TCP garantiza orden y entrega, pero agrega copias y cambios de contexto en cada página. Una implementación sobre RDMA o sobre sockets con memoria compartida reduciría la latencia de los fallos y liberaría CPU para el cálculo. El diseño de mensajes ya separa la cabecera de la carga y valida longitudes, de modo que un transporte nuevo podría conservar el mismo catálogo y cambiar solo el envío. La compresión de páginas o el uso de páginas grandes son variantes que se apoyan en la misma interfaz.

Un panel web con histórico de métricas complementaría los reportes actuales. El reporte HTML muestra una fotografía del estado; una serie temporal de fallos, latencia, tráfico y franjas reconstruidas permitiría ver tendencias y detectar un nodo que se degrada antes de que caiga. El panel podría leer las mismas consultas de \texttt{dsmctl} y guardar las muestras en un archivo o en una base ligera, sin agregar un servicio obligatorio al clúster.

Las aplicaciones de prueba también pueden crecer. Un procesador de imágenes por etapas, un motor de grafos, un compresor paralelo o una rutina de inferencia caben en el modelo de particiones descrito con direcciones de la vista. Cada una ejercita un patrón de acceso distinto y ayuda a medir el costo de la coherencia en cargas reales.
